\documentclass{article}
\usepackage{graphicx}
\usepackage{amsmath}
\usepackage{amssymb}
\usepackage{hyperref}
\usepackage{booktabs}
\usepackage{enumitem}

\title{Joint-Ratio GRPO for Hybrid Diffusion Language Models:\\
Overnight Experimental Report}
\author{Thomas Kidane (setup) \quad Fable 5 (execution)}
\date{July 11, 2026}

\begin{document}
\maketitle

\section{Executive Summary}

\textbf{Verified findings:}
\begin{itemize}
\item The joint ELBO estimator from prior agent work is \textbf{correct}: it reproduces identically (total=10,584,237, contrast vs.\ shuffled=707 nats) on the real 339M CANDI checkpoint.
\item The \textbf{differentiable forward path} through the real CANDI model works: gradients flow correctly with norm $\approx 24.7$.
\item \textbf{Coupling on the real model} is confirmed: the factored-vs-joint importance ratio error has mean absolute value 0.07--0.09 nats and max 0.33 nats across all noise levels $t\in\{0.3, 0.5, 0.7, 0.9\}$. The joint and factored ratios are only weakly correlated (0.25--0.51).
\item \textbf{Joint-ratio GRPO is stable}: 200 steps, no NaN, importance ratios stay near 1.0 ($\text{mean}\approx 0.93$--$1.08$, $\text{std}\approx 0.10$--$0.22$), KL bounded in $[-0.11, 0.06]$, reward best 0.732.
\item \textbf{Factored-ratio GRPO collapses}: same setup, ratios explode to 74$\times$, KL diverges to $-245$, reward crashes from 0.578 to 0.053.
\end{itemize}

\textbf{Preliminary findings (require scaled confirmation):}
\begin{itemize}
\item The reward signal is noisy at $G{=}4$ (expected); a larger group size or longer run is needed to confirm monotonic improvement.
\item The v1 run (sequence-level ratio without per-token normalization) reproduced the known instability from the literature (ratio $>90\times$), confirming per-token normalization is essential.
\end{itemize}

\textbf{Negative results:}
\begin{itemize}
\item The initial coupling measurement in raw 50,258-dimensional one-hot space was numerically catastrophic at low $t$ (values $\sim 10^9$). The correct measurement uses the coordinate-decomposed Gaussian bridge formula.
\item The first GRPO attempt (v1) with sequence-level sum ratio failed identically to the factored baseline---instability is not a joint-vs-factored issue alone, but also a normalization issue. The contribution is that joint+per-token-normalized is the \textit{only} combination that works.
\end{itemize}

\section{Methods}

\subsection{Joint ELBO Estimator}

For the CANDI reverse kernel at a masked position, the discrete channel gives token $y$ with probability $\pi_{\theta,t}(y)$ and the continuous channel follows a posterior-weighted Gaussian bridge:
\begin{equation}
M_\theta(x_s \mid x_t) = \sum_y \pi_{\theta,t}(y)\,\phi_{s,t,y}(x_s \mid x_t)
\end{equation}
where $\phi_{s,t,y}(x_s\mid x_t) = \mathcal{N}(x_s;\, (1{-}a)e_y + a\,x_t,\; v\,I)$ with $a = \sigma_s^2/\sigma_t^2$ and $v = \sigma_s^2(1-\sigma_s^2/\sigma_t^2)$.

The discrete ELBO uses Gauss-Legendre quadrature (3 nodes) over the noise schedule:
\begin{equation}
\text{ELBO}_{\text{disc}} = \sum_{i=1}^{n_\text{quad}} w_i \sum_{b,l} \log p_\theta(x_0^{(b,l)} \mid x_t, t_i)\,\frac{d\alpha}{1-\alpha}\bigg|_{t_i} + \text{prior}
\end{equation}

\subsection{Importance Ratio}

The joint per-step importance ratio for the masked branch is:
\begin{equation}
\rho_{\theta/b}(x_s \mid x_t, z_t{=}m) = \frac{\sum_y \pi_\theta(y)\,\phi_y(x_s\mid x_t)}{\sum_y \pi_b(y)\,\phi_y(x_s\mid x_t)}
\end{equation}
A factored ratio $\pi_\theta(k)/\pi_b(k)$ ignores the continuous observation $x_s$ and the mixture weighting. The \textit{error} of the factored approximation is the change in coupling between old and new policies.

\subsection{GRPO Configuration}

\begin{tabular}{ll}
\toprule
Parameter & Value \\
\midrule
Model & CANDI 339M (GPT-2 tokenizer, OpenWebText) \\
Sequence length & 64 tokens \\
Generation NFE & 16 steps \\
Group size $G$ & 4 \\
Prompts per step & 2 \\
Learning rate & $3\times10^{-6}$ \\
Clip $\epsilon$ & 0.2 \\
KL coefficient & 0.05 \\
Gradient clip & 0.5 \\
Ratio cap & $\exp(\pm2)$ \\
Reward & Fraction of tokens in top-1000 GPT-2 vocabulary \\
\bottomrule
\end{tabular}

The sequence-level log-ratio is the \textit{mean} (not sum) of per-token log-ratios. This normalization is critical for stability.

\section{Results}

\subsection{Gate 1: Estimator Verification}

The joint ELBO estimator was re-verified on the real checkpoint:
\begin{itemize}
\item Total ELBO (real text, $n_\text{quad}$=3): 10,584,237
\item Shuffled contrast: 10,583,530 (difference = 707 nats, positive)
\item Discrete term matches model training loss to relative error $1.5\times10^{-7}$
\item Convergence: $|n_3 - n_5|/|n_5| = 1.08\times10^{-5}$
\end{itemize}

\subsection{Coupling on the Real Checkpoint}

We measure the factored-ratio error by simulating a small policy perturbation (0.1 standard deviations in logit space) and comparing joint vs.\ factored importance ratios:

\begin{table}[h]
\centering
\begin{tabular}{cccccc}
\toprule
$t$ & $N$ & Mean $|\Delta|$ & Max $|\Delta|$ & P90 $|\Delta|$ & Corr(joint, fact) \\
\midrule
0.3 & 74 & 0.089 & 0.313 & 0.167 & --- \\
0.5 & 100 & 0.076 & 0.281 & 0.156 & --- \\
0.7 & 100 & 0.069 & 0.209 & 0.155 & 0.508 \\
0.9 & 100 & 0.076 & 0.334 & 0.177 & 0.248 \\
\bottomrule
\end{tabular}
\caption{Factored-vs-joint ratio error $\Delta = \log r_\text{joint} - \log r_\text{factored}$ on the real 339M CANDI model. The error is substantial and consistent across noise levels.}
\end{table}

\subsection{Joint vs.\ Factored GRPO Ablation}

\begin{table}[h]
\centering
\begin{tabular}{lccccc}
\toprule
Method & Reward$_\text{end}$ & Best & Ratio mean & KL range & Collapse? \\
\midrule
\textbf{Joint (ours)} & 0.588 & \textbf{0.732} & 0.93--1.08 & $[-0.11, 0.06]$ & No \\
Factored & 0.416 & 0.723 & 0.007--74 & $[-245, 0]$ & \textbf{Yes} \\
\bottomrule
\end{tabular}
\caption{200-step GRPO comparison. Same model, reward, hyperparameters. The factored ratio destabilizes within 100 steps while the joint ratio remains controlled throughout.}
\end{table}

Key observations:
\begin{enumerate}
\item The factored method achieves early reward peaks (0.72) before collapsing, because the biased high-variance ratios occasionally align with high-advantage samples.
\item The joint method maintains stability: all ratios stay in $[0.14, 7.4]$ (capped), with most in $[0.8, 1.2]$.
\item The factored method's KL diverges to $-245$, indicating the policy has drifted far from the reference in an uncontrolled direction. The joint method's KL stays bounded.
\end{enumerate}

\section{Threats to Validity}

\begin{itemize}
\item \textbf{Small scale.} $G{=}4$, $L{=}64$, 200 steps. The reward trajectory is noisy and does not show monotonic improvement. A proper run needs $G\geq16$, $L{=}1024$, and $>1000$ steps.
\item \textbf{Single seed.} All results are from seed 42. Seed sensitivity is unknown.
\item \textbf{Toy reward.} ``Fraction of tokens in top-1000'' is verifiable but does not require real language quality. A meaningful evaluation needs a language-model reward or verifier.
\item \textbf{Single time point.} The per-token log-ratio uses $t{=}0.5$ only. A proper estimator should use quadrature over multiple time points (GDPO-style).
\item \textbf{Coupling magnitude.} The 0.07--0.09 nat mean error on the ratio corresponds to a $\sim7\%$--$9\%$ bias on the importance weight. Whether this is catastrophic depends on the advantage distribution and clipping behavior, which we observe empirically (factored collapses, joint does not) but have not characterized analytically on the real model.
\item \textbf{v1 instability.} The first joint run (sequence-level sum ratio) also failed. The working configuration requires BOTH the joint treatment AND per-token normalization. We cannot claim the joint ratio alone is sufficient.
\end{itemize}

\section{Artifacts}

All artifacts are in \texttt{/home/ec2-user/together/shared/artifacts/} and reproducible:

\begin{tabular}{lp{8cm}}
\toprule
File & Description \\
\midrule
\texttt{gate1\_coupling\_fable.json} & Gate 1 verification + initial coupling attempt \\
\texttt{coupling\_real\_checkpoint.json} & Proper coupling measurement (coordinate-decomposed) \\
\texttt{grpo\_real\_run.json} & v1 joint GRPO (unstable, 200 steps) \\
\texttt{grpo\_joint\_v2.json} & v2 joint GRPO (stable, 200 steps) \\
\texttt{grpo\_factored\_run.json} & Factored baseline (collapses, 200 steps) \\
\texttt{grpo\_real\_v2.py} & Joint GRPO implementation \\
\texttt{grpo\_factored.py} & Factored GRPO implementation \\
\texttt{coupling\_v2.py} & Coupling measurement script \\
\texttt{gate1\_coupling.py} & Gate 1 + initial coupling script \\
\bottomrule
\end{tabular}

To regenerate:
\begin{verbatim}
cd /home/ec2-user/together/ashir/candi-diffusion
PYTHONPATH=/home/ec2-user/trajectoryforge/venv/lib64/python3.9/site-packages:.
python3 /home/ec2-user/together/shared/artifacts/gate1_coupling.py
python3 /home/ec2-user/together/shared/artifacts/coupling_v2.py
python3 /home/ec2-user/together/shared/artifacts/grpo_real_v2.py
python3 /home/ec2-user/together/shared/artifacts/grpo_factored.py
\end{verbatim}

\section{Conclusions}

The joint-ratio formulation for GRPO on CANDI is both theoretically motivated and empirically necessary. The factored ratio introduces a bias of 0.07--0.09 nats (7--9\% on the importance weight) that is sufficient to destabilize training. With proper per-token normalization and ratio capping, the joint GRPO loop runs stably on the real 339M-parameter checkpoint and produces a reward signal.

The immediate next steps for a submission-quality result are: (1) scale to $G\geq16$ and $L{=}1024$ with a meaningful reward; (2) run for $>$1000 steps to show monotonic improvement; (3) compare against GDPO-style variance-decomposed estimator extended to the joint case.

\end{document}
